\section{A multi-level evaluation protocol}
\label{sec:eval-suite}

\F{F1} showed that surface-proximity metrics can reward an anatomically worse
fit. A direct adversarial test sharpens this: deliberately snapping distal
vertices onto the nearest target surface improves a per-part distance metric
by $53\%$ while edge-length distortion rises $7.5$-fold. A metric that a
targeted attack can improve by damaging the mesh cannot be the sole judge of
a fit. The protocol therefore has three evidence levels, always used
together:

\begin{center}
\begin{tabular}{@{}p{0.18\linewidth}p{0.34\linewidth}p{0.36\linewidth}@{}}
\toprule
\textbf{Level} & \textbf{Metrics} & \textbf{Question answered} \\
\midrule
1. Surface proximity & Chamfer, F-score, Hausdorff, per-part variants &
  Does the fitted mesh cover the target? \\
2. Target-independent integrity & Edge-length distortion, Laplacian /
  as-rigid-as-possible deviation~\cite{arap}, normal consistency &
  Was that coverage bought by damaging the mesh? \\
3. Anatomical ground truth & Synthetic vertex-identity error
  (\S\ref{sec:synthetic-instrument}); expert joint error on the corrected
  benchmark (\S\ref{sec:skeleton-underdetermined}) &
  Is vertex $i$ on the anatomical locus it represents? \\
\bottomrule
\end{tabular}
\end{center}

An intervention is adopted only if it does not degrade levels 1--2 and is
checked at level 3 wherever ground truth exists.
